\documentclass[11pt]{article}

\usepackage[T1]{fontenc}
\usepackage{lmodern}
\usepackage{microtype}
\usepackage{amsmath,amssymb,amsthm}
\usepackage{booktabs}
\usepackage{enumitem}
\usepackage{xcolor}
\usepackage[margin=1in]{geometry}
\usepackage{xurl}
\usepackage[
  colorlinks=true,
  linkcolor=blue!55!black,
  citecolor=blue!55!black,
  urlcolor=blue!55!black,
  pdfauthor={Austin Lin Gibbons},
  pdftitle={Nonmonotone Rowley Templates and the Bound R5(5) >= 44338}
]{hyperref}
\usepackage{fancyhdr}

\setlist[itemize]{leftmargin=1.5em,itemsep=2pt,topsep=4pt}
\setlist[enumerate]{leftmargin=1.7em,itemsep=2pt,topsep=4pt}
\setlength{\parindent}{0pt}
\setlength{\parskip}{5pt plus 1pt minus 1pt}
\clubpenalty=10000
\widowpenalty=10000
\displaywidowpenalty=10000

\pagestyle{fancy}
\fancyhf{}
\fancyhead[L]{\small Nonmonotone Rowley templates}
\fancyhead[R]{\small\thepage}
\renewcommand{\headrulewidth}{0.3pt}

\newtheorem{theorem}{Theorem}[section]
\newtheorem{proposition}[theorem]{Proposition}
\newtheorem{lemma}[theorem]{Lemma}
\newtheorem{corollary}[theorem]{Corollary}
\theoremstyle{definition}
\newtheorem{definition}[theorem]{Definition}
\newtheorem{certificate}[theorem]{Certificate}
\theoremstyle{remark}
\newtheorem{remark}[theorem]{Remark}

\newcommand{\Z}{\mathbb Z}
\newcommand{\Rfive}{R_5(5)}
\newcommand{\code}[1]{\texttt{#1}}
\newcommand{\sha}[1]{\begingroup\ttfamily\scriptsize\urlstyle{same}\nolinkurl{#1}\endgroup}
\newcommand{\hashbreak}[2]{{\ttfamily\scriptsize #1\linebreak #2}}

\title{\bfseries Nonmonotone Rowley Templates and the Bound\\[3pt]
  \texorpdfstring{\boldmath$\Rfive\ge44{,}338$}{R5(5) >= 44338}}
\author{Austin Lin Gibbons\thanks{Draft for independent review. Correctness,
priority, and publication readiness are treated as separate questions.}}
\date{23 July 2026}

\begin{document}
\maketitle

\begin{abstract}
We give an explicit effective Rowley $(5,5,3)$ template of order $99$ with
extension parameter $\phi=40$. Its inherited distance colours remain
$K_5$-free under period-$98$ repetition through the exact cutoff
$4(99-2)=388$, and its third distance class is sum-free. Combining it with
an explicit linear $(5,5,5;453)$ prototype by Rowley's generalized
construction yields a five-colouring of $K_{44,337}$ without a
monochromatic $K_5$. Consequently,
\[
  R_5(5)\ge44{,}338.
\]
We also certify that no effective $(5,5,3)$ template of order $98$ has
$\phi\ge40$. Since templates with $\phi=40$ exist at orders $97$ and $99$,
feasibility at fixed extension threshold is nonmonotone in the order.
The negative result is exhaustive: one DRAT certificate proves that the
first forty distances have only twenty-two admissible two-colour prefixes,
and eleven further DRAT certificates exclude every unrestricted tail up to
global colour exchange. An independent semantic checker reconstructs all
CNFs from explicit sum and five-vertex witnesses before proof checking.
The positive objects are checked by independent Python and C++ clique
searches, and the full compound word is reconstructed independently.
\end{abstract}

\medskip
\noindent\textbf{Keywords.} Multicolour Ramsey numbers; linear colourings;
triangle-free templates; SAT certificates; DRAT; nonmonotonicity.

\noindent\textbf{MSC 2020.} 05C55, 05D10, 68V15.

\section{Introduction}

For positive integers $r$ and $k$, the diagonal multicolour Ramsey number
$R_r(k)$ is the least $N$ such that every $r$-edge-colouring of $K_N$
contains a monochromatic $K_k$. Thus an explicit $r$-colouring of
$K_M$ without a monochromatic $K_k$ proves $R_r(k)\ge M+1$.

Rowley developed a composition method in which a linear colouring with a
distinguished triangle-free template is combined with a second linear
Ramsey colouring \cite{rowley-general,rowley-sat}. His published effective
$(5,5,3)$ template of order $93$ and parameter $\phi=40$, combined with an
order-$453$ prototype, gives
\[
  R_5(5)\ge41{,}626.
\]
Revision 18 of the standard survey records the same value
\cite{radziszowski}.

This note provides an order-$99$ template at the same value of $\phi$.
The compound order is
\[
  (99-1)(453-1)+1+40=44{,}337,
\]
giving the asserted lower bound $44{,}338$. The complete $44{,}336$-symbol
distance word is included in the accompanying artifact package.

The intervening failure at order $98$ is not merely a failed search.
It is proved by exhaustive case decomposition and checkable
unsatisfiability certificates. Together with valid templates at orders $97$
and $99$, this proves that the set of feasible orders at fixed $\phi$ need
not be an interval. The order-$99$ word uses a non-cyclic forty-symbol
prefix, making symmetry breaking part of the mathematical mechanism rather
than a search convenience.

The latest survey checked here is revision 18, dated 24 April 2026
\cite{radziszowski}. Bounded searches found no competing order-$99$
template or the numerical bound $44{,}338$. This is not a proof of priority;
the claims of this note are correctness claims pending external
current-literature review.

\section{Linear templates}

\begin{definition}[Linear distance colouring]
Let $m\ge2$. A linear three-colouring of $K_m$ is a word
\[
  c:\{1,\ldots,m-1\}\longrightarrow\{1,2,3\},
\]
where edge $\{i,j\}$, $i<j$, receives colour $c(j-i)$. Put
$p=m-1$ and
\[
  A_1=c^{-1}(1),\qquad A_2=c^{-1}(2),\qquad T=c^{-1}(3).
\]
\end{definition}

The set $T$ is the distinguished template colour. Its linear graph is
triangle-free exactly when
\begin{equation}
  \nexists\,a,b>0\quad a,b,a+b\in T,
  \label{eq:sumfree}
\end{equation}
including the case $a=b$. The extension parameter is
\[
  \phi=\min(T)-1.
\]

For $s\in\{1,2\}$, repeat inherited colour $s$ with period $p$ by assigning
positive distance $d$ the base colour
\begin{equation}
  c\!\left(((d-1)\bmod p)+1\right).
  \label{eq:periodic}
\end{equation}
The representative of residue zero is therefore $p$, not $0$.

\begin{definition}[Effective $(5,5,3)$ template]
An order-$m$ word is effective with parameter $\phi$ if:
\begin{enumerate}
  \item $p=m-1$ belongs to $T$ and $\min(T)-1=\phi$;
  \item $T$ satisfies \eqref{eq:sumfree}; and
  \item for $s=1,2$, the periodically repeated colour in
  \eqref{eq:periodic} contains no $K_5$.
\end{enumerate}
\end{definition}

The infinite-looking final condition has an exact finite form.

\begin{lemma}[Directed-Cayley cutoff]
\label{lem:cutoff}
Fix an inherited colour $s$ and period $p$. Define a directed relation on
$\Z_p$ by
\[
  x\longrightarrow y
  \quad\Longleftrightarrow\quad
  c\!\left(((y-x-1)\bmod p)+1\right)=s.
\]
The periodic colour contains a $K_5$ if and only if there are distinct
$r_0,\ldots,r_4\in\Z_p$ such that
\[
  r_i\longrightarrow r_j\qquad(0\le i<j\le4).
\]
If such a witness exists, one exists on integer vertices contained in
$\{0,\ldots,4(p-1)\}$.
\end{lemma}

\begin{proof}
An increasing integer witness reduces modulo $p$ to the displayed ordered
relation. Equal residues cannot occur, because their difference is a
multiple of $p$, represented by $p\in T$ rather than inherited colour $s$.

Conversely, for $i=0,\ldots,3$, let $g_i\in\{1,\ldots,p-1\}$ be the least
positive representative of $r_{i+1}-r_i$. Put $x_0=0$ and
$x_{i+1}=x_i+g_i$. Then $x_j-x_i$ has residue $r_j-r_i$ for every
$i<j$, so the five $x_i$ form a monochromatic clique. Moreover
\[
  x_4=g_0+g_1+g_2+g_3\le4(p-1).
\]
\end{proof}

Lemma~\ref{lem:cutoff} is the $(5,5,3)$ specialization of the compression
mechanism in Rowley's Theorem 3.2 \cite{rowley-general}. It also identifies
the native finite obstruction: an ordered transitive $K_5$ in a directed
Cayley relation.

\section{The order-99 certificate}

\begin{certificate}[Order-$99$ word]
\label{cert:order99}
Concatenate the following two lines:
\begin{quote}
\ttfamily\footnotesize
2111211211122211222222221122211121121112321132111133232323331333331333\\
3113333133323332333111231123
\end{quote}
The resulting word has length $98$.
\end{certificate}

\begin{proposition}
\label{prop:order99}
Certificate~\ref{cert:order99} is an effective $(5,5,3)$ template of order
$99$ with $\phi=40$.
\end{proposition}

\begin{proof}
Direct enumeration gives $98\in T$, $\min(T)=41$, $|T|=33$, and no
$a,b,a+b\in T$. Two independent exact clique searches form the repeated
inherited-colour graphs on vertices $0,\ldots,388$ and reject $K_5$ in
both colours. By Lemma~\ref{lem:cutoff}, span $388=4(98-1)$ is complete.
The Python implementation visits $1{,}645{,}997$ and $917{,}897$ search
nodes for the two colours; the independent C++ implementation also returns
\code{VALID}.
\end{proof}

The canonical file is
\code{results/order99\_linear\_prefix8.template}, with SHA-256
\begin{center}
\sha{2643001cc425898d584bd374e20928d1dbc6a72fd4011711151343d0ad072966}.
\end{center}
Its first forty symbols contain twenty occurrences of each inherited colour
and are not cyclically reflected.

\section{The order-453 prototype}

The second input is the $452$-symbol word in Appendix~\ref{app:prototype}.
It is recovered from the archived order-$977$ template in Rowley's ancillary
data: the first $452$ distances are the order-$453$ prototype and distance
$453$ begins the newly introduced template colour \cite{rowley-sat}.

\begin{proposition}
\label{prop:prototype}
The word in Appendix~\ref{app:prototype} defines a cyclic linear
three-colouring of $K_{453}$ with no monochromatic $K_5$.
\end{proposition}

\begin{proof}
Independent Python bit-mask and C++ ordered-vector clique searches
exhaustively reject $K_5$ in each colour. Their exact node counts are,
respectively,
\[
\begin{array}{c|rrr}
 &1&2&3\\ \hline
\text{Python}&1{,}040{,}265&5{,}811{,}421&782{,}585\\
\text{C++}&237{,}855&837{,}235&201{,}199.
\end{array}
\]
Direct comparison also gives $v(d)=v(453-d)$ for
$1\le d\le452$.
\end{proof}

The artifact
\code{sources/rowley\_exoo\_order453.prototype} has SHA-256
\begin{center}
\sha{19c97e6279c184f6f462786cadda4b7c9773d870a5b680a04eb1503ef384a2d0}.
\end{center}
Once the word is explicit and Proposition~\ref{prop:prototype} is checked,
its historical recovery is no longer a logical dependency.

\section{The compound colouring}

Set $p=98$. Let $A_1,A_2,T$ be the classes of
Certificate~\ref{cert:order99}, and let $B_1,B_2,B_3$ be the prototype
classes. For $s=1,2$, define
\begin{equation}
\begin{split}
A'_s={}&
\{\ell+(\mu-1)p:\ell\in A_s,\ 1\le\mu\le452\}\\
&{}\cup
\{\ell+452p:\ell\in A_s,\ 1\le\ell\le40\}.
\end{split}
\label{eq:inherited}
\end{equation}
For $j=1,2,3$, put
\begin{equation}
  B'_j=\{\ell+(\mu-1)p:\ell\in T,\ \mu\in B_j\}.
  \label{eq:prototype}
\end{equation}

\begin{lemma}
\label{lem:partition}
The five sets $A'_1,A'_2,B'_1,B'_2,B'_3$ partition
$\{1,\ldots,44{,}336\}$.
\end{lemma}

\begin{proof}
For $1\le d\le452p$, write uniquely
$d=\ell+(\mu-1)p$ with $1\le\ell\le p$ and
$1\le\mu\le452$. If $\ell\in A_s$, equation \eqref{eq:inherited} assigns
$d$ to $A'_s$. If $\ell\in T$, exactly one $B_j$ contains $\mu$, and
\eqref{eq:prototype} assigns $d$ to $B'_j$.

The remaining distances are $452p+\ell$ with $1\le\ell\le40$. Since
$T\cap[1,40]=\varnothing$, the second line of
\eqref{eq:inherited} assigns each uniquely.
\end{proof}

The five class sizes are
\[
\begin{array}{c|rrrrr|r}
\text{class}&A'_1&A'_2&B'_1&B'_2&B'_3&\text{total}\\ \hline
\text{size}&16{,}292&13{,}128&4{,}488&6{,}402&4{,}026&44{,}336.
\end{array}
\]

\begin{theorem}
\label{thm:main}
The distance classes in \eqref{eq:inherited}--\eqref{eq:prototype} define a
five-colouring of $K_{44,337}$ without a monochromatic $K_5$. Consequently,
\[
  \Rfive\ge44{,}338.
\]
\end{theorem}

\begin{proof}
Lemma~\ref{lem:partition} defines the colouring. We exclude the two types of
output colour separately.

\emph{Inherited colours.}
Suppose output colour $s\in\{1,2\}$ contains a $K_5$. Translate its least
vertex to zero and reduce all five vertices modulo $p$. Every ordered
pairwise residue difference lies in $A_s$. Lemma~\ref{lem:cutoff} therefore
produces a $K_5$ in the finite repeated graph on
$0,\ldots,388$, contradicting Proposition~\ref{prop:order99}.

\emph{Prototype-derived colours.}
Fix $j\in\{1,2,3\}$ and suppose output colour $j+2$ contains a $K_5$.
Translate its least vertex to zero. Each other vertex has a unique
representation
\[
  d_i=t_i+(\mu_i-1)p,\qquad t_i\in T,\quad\mu_i\in B_j.
\]
For $d_i<d_k$, membership of $d_k-d_i$ in $B'_j$ forces
$t_i\ge t_k$: if $t_i<t_k$, then the positive residue
$t_k-t_i$ would belong to $T$, contradicting
$t_i+(t_k-t_i)=t_k$ and sum-freeness.

Thus the residues $t_i$ are nonincreasing along the ordered vertices. The
prototype index of $d_k-d_i$ is then $\mu_k-\mu_i$, whether
$t_i=t_k$ or $t_i>t_k$. Hence every positive difference
$\mu_k-\mu_i$ lies in $B_j$. The indices are strictly increasing, so the
five indices, including zero for the translated least vertex, form a
monochromatic $K_5$ in prototype colour $j$, contradicting
Proposition~\ref{prop:prototype}.

All five colours are $K_5$-free. The order calculation gives
\[
  98\cdot452+41=44{,}337,
\]
and the Ramsey inequality follows.
\end{proof}

The complete word is
\code{results/r5\_5\_order44337.linear-coloring}, with SHA-256
\begin{center}
\sha{274acbf17bf7732b16ef7d20c97486eb469486907fd1357c16990ed4332f7158}.
\end{center}
It is linear and non-cyclic. A generator implementing the block formula and
an independent literal set-union reconstruction agree on all $44{,}336$
symbols.

\section{Certified nonexistence at order 98}

\begin{theorem}
\label{thm:hole}
There is no effective Rowley $(5,5,3)$ template of order $98$ with
$\phi\ge40$.
\end{theorem}

\begin{proof}[Certificate architecture]
Any candidate has period $97$, terminal distance $97$ in $T$, and
$T\cap[1,40]=\varnothing$. Thus the first forty distances form a two-colour
linear colouring of $K_{41}$ with neither colour containing $K_5$.

\emph{Stage 1: exhaustive prefixes.}
For each five-set $V\subseteq\{0,\ldots,40\}$, let
\[
  \Delta(V)=\{|u-v|:u,v\in V,\ u\ne v\}.
\]
If Boolean variable $x_d$ selects inherited colour 1 at distance $d$, the
two clauses
\[
  \bigvee_{d\in\Delta(V)}x_d,\qquad
  \bigvee_{d\in\Delta(V)}\neg x_d
\]
exclude a monochromatic $K_5$ in colours 2 and 1. Deduplication over all
$V$ gives $45{,}374$ distinct distance sets and $90{,}748$ clauses.

There are 22 satisfying prefix words, stored as 11 representatives up to
global exchange of colours 1 and 2. Adding blocking clauses for those 22
words makes the prefix CNF unsatisfiable. The file
\code{prefix41\_exhaust.drat} is a verified DRAT proof of that UNSAT
instance. Independently, the semantic checker regenerates all $45{,}374$
sets, checks each representative directly, reconstructs the CNF
byte-for-byte, and invokes the proof checker.

\emph{Stage 2: unrestricted tails.}
For each of the 11 representatives, introduce three one-hot variables for
every distance $1,\ldots,97$. Fix only the selected first forty symbols and
the terminal template colour. Every distance $41,\ldots,96$ is allowed all
three colours.

For every $a,b>0$ with $a+b\le97$, include the clause forbidding all of
$a,b,a+b$ from lying in $T$. Each inherited-colour obstruction clause
carries an explicit ordered witness
\[
  0=v_0<v_1<\cdots<v_4\le384.
\]
Its period-$97$ distance coordinates are
\[
  D(V)=
  \left\{((v_j-v_i-1)\bmod97)+1:0\le i<j\le4\right\},
\]
and the clause forbids all coordinates of $D(V)$ from receiving the
corresponding inherited colour. Every such clause is therefore a valid
necessary condition for a template.

All 11 tail CNFs are UNSAT, with verified DRAT proofs. The semantic checker
validates every sum identity and every five-vertex witness, then reconstructs
each CNF byte-for-byte before proof checking. Exchanging inherited colours
covers the complementary 11 prefixes. Hence every possible first forty
symbols and every unrestricted tail are eliminated.
\end{proof}

\begin{remark}
The tail formulas need not contain every possible $K_5$ clause. Each recorded
clause is valid, and UNSAT of a relaxation containing only recorded valid
obstructions already proves that no candidate exists.
\end{remark}

The certificate packet is
\code{certificates/order98\_phi40\_exhaustion/}. It contains 12 CNF/DRAT
pairs, semantic clause sidecars, prefix representatives, hash manifests, and
the independent checker. The checker compiles against a pinned
MIT-licensed copy of \code{drat-trim} \cite{drat-trim}.

\begin{corollary}[Nonmonotonicity]
\label{cor:nonmonotone}
At fixed threshold $\phi=40$, existence of an effective $(5,5,3)$ template
is not monotone in the template order.
\end{corollary}

\begin{proof}
The frozen word \code{results/order97\_reflected.template} is accepted by
both exact verifiers, Theorem~\ref{thm:hole} excludes order $98$, and
Proposition~\ref{prop:order99} supplies order $99$.
\end{proof}

\section{Verification and reproducibility}

All release-critical checks run with Python 3, a C compiler, a C++17
compiler, and a POSIX shell:
\begin{quote}
\ttfamily\small
sh tests/run\_certificate\_checks.sh
\end{quote}
The suite has no third-party Python dependency. It:
\begin{enumerate}
  \item runs positive and deliberately corrupted fixtures against independent
  template verifiers;
  \item checks all three prototype colours with two clique algorithms;
  \item checks the order-$97$ and order-$99$ words with both algorithms;
  \item regenerates the complete compound and independently reconstructs all
  output distance sets;
  \item compiles the pinned \code{drat-trim.c}; and
  \item semantically reconstructs all order-$98$ CNFs and verifies all 12
  DRAT proofs.
\end{enumerate}
GitHub Actions executes the same command on Ubuntu and macOS.

The two clique implementations share only the mathematical specification.
The Python code uses arbitrary-precision bit masks; the C++ code uses
ordered candidate vectors. Both exhaustively recurse on a candidate set
intersected with the forward neighbours of each selected vertex. A branch is
pruned only when fewer candidates remain than the number of required clique
vertices.

The $44{,}337$-vertex graph is not checked by a naive all-vertex clique
enumeration. Its correctness follows from Theorem~\ref{thm:main}, whose two
finite inputs are checked exactly.

\section{Scope and literature status}

Theorem~\ref{thm:hole} is exact only at order $98$ and threshold
$\phi\ge40$. It says nothing about smaller $\phi$ or later orders. Solver
reconnaissance at orders $100$ through $105$ produced no positive object,
but those runs are not certificate-backed and are not claims.

Rowley's 2022 paper gives $R_5(5)\ge41{,}626$
\cite{rowley-sat}; revision 18 of the dynamic survey records that value
\cite{radziszowski}. The asymptotic update of Campos and Pohoata does not
state a competing finite value \cite{campos-pohoata}. Searches for
$44{,}338$, an order-$99$ effective template, and follow-ups to Rowley's
construction found no collision in the checked corpus.

Accordingly, the present status is:
\begin{center}
\textbf{finite correctness certified locally; external priority review
required.}
\end{center}

\appendix

\section{Order-453 prototype word}
\label{app:prototype}

Concatenate the following seven lines exactly:
\begin{quote}
\ttfamily\scriptsize
1222333322121112131212231312221323223233222132322323322111112212111221\\
3132231212232313223121323231222132332111122233231221212233323221121321\\
2111223121122233232332332323322233233211111221111123131222213122313122\\
1213222312112223322211213222312122131322131222213132111112211111233233\\
2223323233233232332221121322111212312112232333221212213233222111123323\\
1222132323121322313232212132231312211121221111122332322323122233232232\\
31222131322121312111212233332221
\end{quote}

\section{Principal artifact hashes}

\begin{center}
\begin{tabular}{@{}lp{0.58\textwidth}@{}}
\toprule
Artifact & SHA-256\\
\midrule
Order-$99$ template &
\hashbreak{2643001cc425898d584bd374e20928d1}
{dbc6a72fd4011711151343d0ad072966}\\
Order-$97$ template &
\hashbreak{3287e4737fed6d355d59c4c08c3140a}
{829508039293914109c78000c7ba3a309}\\
Order-$453$ prototype &
\hashbreak{19c97e6279c184f6f462786cadda4b7c}
{9773d870a5b680a04eb1503ef384a2d0}\\
Complete compound &
\hashbreak{274acbf17bf7732b16ef7d20c97486eb}
{469486907fd1357c16990ed4332f7158}\\
Prefix-exhaustion CNF &
\hashbreak{aec46291803dae6b29e180bded569d24}
{102fa5234aa3b8be096177b3eac1c179}\\
Prefix-exhaustion DRAT &
\hashbreak{02ee6ff73c60629c7b9efbac2a95d32}
{6537d23c22e4e83a9380bb7c392aef57d}\\
\bottomrule
\end{tabular}
\end{center}

\section*{Acknowledgement}

Computational exploration and technical drafting used automated coding and
language-model assistance. Every mathematical claim in this note is tied to
explicit source text, finite objects, or independently checkable
certificates; no model output is used as authority.

\begin{thebibliography}{9}
\raggedright

\bibitem{rowley-general}
F. Rowley,
\newblock A generalised linear Ramsey graph construction,
\newblock \emph{Australasian Journal of Combinatorics} \textbf{81}(2)
(2021), 245--256.
\newblock \url{https://ajc.maths.uq.edu.au/pdf/81/ajc_v81_p245.pdf}.

\bibitem{rowley-sat}
F. Rowley,
\newblock Improved lower bounds for multicolour Ramsey numbers using
SAT-solvers,
\newblock \emph{arXiv:2203.13476v3}, 2022.
\newblock \url{https://arxiv.org/abs/2203.13476}.

\bibitem{radziszowski}
S. P. Radziszowski,
\newblock Small Ramsey numbers,
\newblock \emph{Electronic Journal of Combinatorics}, Dynamic Survey DS1,
revision 18, 24 April 2026.
\newblock \url{https://doi.org/10.37236/21}.

\bibitem{campos-pohoata}
M. Campos and C. Pohoata,
\newblock An update on multicolor Ramsey lower bounds,
\newblock \emph{arXiv:2601.15183}, 2026.
\newblock \url{https://arxiv.org/abs/2601.15183}.

\bibitem{drat-trim}
M. Heule and N. Wetzler,
\newblock DRAT-trim proof checker,
\newblock \url{https://github.com/marijnheule/drat-trim}.

\end{thebibliography}

\end{document}
